\begin{textsample}{NEG Dim 5 – vem\_ed\_02 – Score 1.00 – vem\_ed\_02/t005.txt}  \label{ex:f5_neg_088}
PROCURA-SE Quer desenvolver um Projeto Colaborativo Internacional (PCI)?
Parceiros em Instituições de Ensino estrangeiras buscam pares nas Fatecs em diversas áreas, como: Competências interculturais Gerentes internacionais Mineração e geoprocessamento Mapa de riscos Saúde Meio Ambiente e Sustentabilidade Artes Álgebra e Cálculo História e Ciência Política Caso tenha interesse, basta \textbf{escrever} para cesu.pci@cps.sp.gov.br e juntar- se às nossas equipes no Teams.
QUEBRA-GELO Esta segunda edição de VEm com PCI traz uma amostra de esforços de internacionalização nas Fatecs.
Em julho, a equipe dos PCIs organizou uma palestra oferecida por videoconferência pelo professor Michael Fields, da Eastern Oregon University (EUA), para as Fatecs de Ferraz de Vasconcelos, Itaquaquecetuba e Mogi das Cruzes.
A seção “Quem é Quem” apresenta brevemente Sally Crimmins-Villela, da State University of New York (SUNY), líder na área de intercâmbios virtuais.
Ela leu e comentou a primeira edição de VEm com PCI: isso alegra e orgulha muito nosso time!
Projetos colaborativos que amadurecem e perduram por mais de um semestre demonstram o engajamento dos parceiros na evolução de suas pesquisas.
Por isso, neste número, apresentamos dois exemplos de trabalhos que iniciaram em 2019 com instituições norte-americanas e prosseguem neste ano: a campanha de conscientização contra violência de gênero, elaborada pela Fatec Barueri com a Florida International University (FIU); e o estudo que Fatec Bauru e BridgeValley Community \& Technical College (EUA) elaboram sobre manutenção de equipamentos médicos (neste segundo semestre de 2020, as pesquisas irão focar nos protocolos de higienização para o combate ao coronavírus).
Boa leitura!

% matched lemmas: escrever
\end{textsample}
